\chapter{Introduction}
\label{ch:introduction}

\section{Introduction}
\label{sec:introduction}

Software development is changing quickly. Large language models (LLMs) first
entered developers' work as code-completion assistants, such as GitHub Copilot
and ChatGPT, that suggest lines or functions inside the editor
\cite{yu2024where}. A newer generation of tools, called \emph{AI coding agents},
goes further. Agents such as OpenAI Codex, GitHub Copilot's coding agent, Devin,
Cursor, Claude Code, and Google Jules receive a task, edit several files, run
commands, and open a complete pull request (PR) on GitHub without a human writing
the code \cite{li2026aidev,watanabe2025agentic}. Human maintainers then review
these PRs and decide whether to merge them into the project.

The scale of this change is large. The AIDev dataset, which collects PRs opened
by AI coding agents on public GitHub repositories, grew from about 0.9 million
agent PRs in its third version to more than 7.6 million PRs in its fifth version,
covering activity up to March 2026 \cite{li2026aidev,aidev2026hf}. Many of these
PRs are accepted: in a study of PRs created with Claude Code, 83.8\% were
eventually merged \cite{watanabe2025agentic}. AI-generated code is therefore no
longer an experiment; it is becoming a normal part of the code base of real
open-source projects.

Most research on AI-generated code looks at the moment \emph{before} the merge:
whether the code is correct on benchmarks, what errors LLMs make
\cite{wang2025towards}, whether agent PRs are accepted, and what happens during
review \cite{watanabe2025agentic}. Much less is known about what happens
\emph{after} the merge. Merging a PR does not mean that the code is finished.
Once the code is part of the project, other developers or agents may come back to
it to correct a defect, a missing check, a design weakness, or a performance
problem. In software engineering, this kind of work is called \emph{corrective
maintenance}: changes made to remove faults or deficiencies from code that is
already in use \cite{swanson1976dimensions}.

Recent studies give the first signs that merged agent code does need such
corrections. Takerngsaksiri et al. linked merged agent PRs to later PRs in the
same repositories and found that 22.9\% of merged agent PRs were followed by a
candidate fix, and 4.5\% by a verified fix, within 30 days; agent merges were
fixed at 1.62 times the odds of human merges in the same repositories
\cite{takerngsaksiri2026finishes}. They also found that 69.6\% of the verified
fixes to agent merges came from the same agent. Sawada et al. tracked files
created by agents and found that humans authored 83.2\% of the later commits to
those files \cite{sawada2026maintenance}. Khemissi et al. showed that when
agent PRs are referenced later, it is often to integrate or fix them
\cite{khemissi2026humans}. Xia and Miller measured how long agentic code
survives after it is merged \cite{xia2026postmerge}.

These studies answer an important first question: \emph{how often} is agent code
changed after the merge. They leave open the questions that matter most for
improving AI coding agents and the processes around them: \emph{why} the code had
to be corrected, \emph{what} was wrong with it, and \emph{how} it was corrected.
They also rely on data that stops in mid-2025, before the large growth of agent
activity, and they mostly count corrections that are explicitly labelled as bug
fixes. A deficiency in AI code is not always repaired by a PR called ``fix'': it
may be corrected through a refactoring, a performance change, or added error
handling. This thesis addresses these gaps.

\section{Purpose}
\label{sec:purpose}

The purpose of this thesis is to understand how AI-generated code is corrected
after it becomes part of a software project. Concretely, the study follows
merged pull requests written by AI coding agents in open-source Python
projects, finds the later pull requests that change the same code, and then
analyses those follow-up changes to separate real corrections of the AI code from
changes made for other reasons, such as new features.

For each correction, the study aims to explain why it was needed, what
deficiency it addressed in the original AI code, and how the deficiency was
repaired. By comparing these corrections across agents and projects, the thesis
aims to describe recurring weaknesses of agent-written code that are not caught
by review, tests, or continuous integration before the merge.

The thesis is an empirical study. It does not propose a new AI agent or a new
tool. It builds on an existing, published linking method
\cite{takerngsaksiri2026finishes}, adapts it to the most recent version of the
AIDev dataset, and combines automatic mining with manual analysis of the code
changes.

\section{Scope}
\label{sec:scope}

The study includes:

\begin{itemize}
  \item \textbf{AI-generated pull requests only.} The starting point is pull
  requests opened by AI coding agents (OpenAI Codex, GitHub Copilot, Devin,
  Cursor, Claude Code, and Google Jules), as identified by the AIDev dataset
  \cite{li2026aidev}. Code written by humans is not studied as a starting point.
  \item \textbf{Merged pull requests.} Only agent PRs that were merged are
  followed, because only merged code can need maintenance in the project.
  \item \textbf{Python projects.} The study focuses on repositories whose main
  language is Python, so that the code changes can be analysed in depth with one
  set of tools.
  \item \textbf{Popular repositories.} Following prior work, only repositories
  with at least 500 GitHub stars are included, to focus on maintained projects
  rather than personal experiments \cite{takerngsaksiri2026finishes}.
  \item \textbf{Follow-up changes within a limited time window.} A later pull
  request is considered a follow-up when it is merged in the same repository
  within 30 or 60 days after the agent PR and changes at least one of the same
  files.
  \item \textbf{Corrective maintenance.} Among the follow-ups, the study analyses
  the changes that solve a problem in the AI code: a bug, a missing check, a
  robustness or error-handling problem, a performance problem, or a design or
  maintainability deficiency.
\end{itemize}

The study excludes:

\begin{itemize}
  \item follow-up changes that only add new features, update documentation, or
  change style without correcting the AI code;
  \item agent pull requests that were closed without being merged;
  \item code generated by AI outside pull requests (for example, code pasted from
  a chat assistant), which cannot be reliably attributed;
  \item programming languages other than Python;
  \item private or industrial repositories.
\end{itemize}

\section{Assumptions}
\label{sec:assumptions}

The study relies on the following assumptions:

\begin{enumerate}
  \item \textbf{Agent attribution is correct.} A pull request labelled as
  agent-authored in AIDev was generated by that agent. AIDev identifies agent PRs
  from signals such as the bot account, the branch name, or a co-author line
  \cite{li2026aidev}; human edits inside an agent PR before the merge are
  possible and are treated as part of the agent PR.
  \item \textbf{Shared files indicate related work.} When a later PR changes a
  file that the agent PR also changed, shortly after the merge, it is a
  candidate follow-up of the agent PR. This is a heuristic: a shared file does
  not prove that the later change corrects the agent code, which is why the
  candidates are verified by manual analysis.
  \item \textbf{The time window is meaningful.} Corrections that are linked to
  the agent PR usually happen within weeks of the merge. Later changes are
  less likely to be about the original agent code.
  \item \textbf{Public repository history is complete enough.} The pull requests,
  commits, and changed files available in AIDev and on GitHub reflect the
  changes made to the code, even though some data may be missing (for example,
  deleted pull requests or very large PRs).
  \item \textbf{Python projects are representative of agent use.} Python is one
  of the most common languages of agent PRs: about one quarter of the agent PRs
  in AIDev version~4 (687,771 of 2,743,854) are in Python repositories. Findings
  on Python projects therefore describe a large part of real agent activity,
  although they may not generalize to all languages.
\end{enumerate}

\section{Research Question}
\label{sec:research-question}

\subsection{Primary Research Question}
\label{subsec:primary-research-question}

\begin{quote}
\emph{How is code written by AI coding agents corrected after it is merged into
open-source Python projects?}
\end{quote}

\subsection{Secondary Research Questions}
\label{subsec:secondary-research-questions}

\begin{description}
  \item[RQ1 -- Why and how often?] How often, and how soon after the merge, do
  merged agent pull requests receive follow-up changes that correct them, and
  how does this differ between AI coding agents?
  \item[RQ2 -- What was wrong?] What kinds of deficiencies in the agent-written
  code do these corrections address (for example, functional bugs, missing
  checks, error handling, performance, or design problems)?
  \item[RQ3 -- How is it fixed?] How are these deficiencies corrected: what kind
  and size of code change is made, and who makes it (the same agent, another
  agent, or a human developer)?
\end{description}

\subsection{Objectives}
\label{subsec:objectives}

To answer these questions, the thesis has the following objectives:

\begin{enumerate}
  \item Review the empirical studies on AI-generated code and agent pull
  requests, and the datasets they share, to select a suitable data source.
  \item Build a clean dataset of merged agent pull requests in Python
  repositories from the most recent version of AIDev.
  \item Adapt and validate an existing pull-request linker
  \cite{takerngsaksiri2026finishes}, first by reproducing its published results
  on the original data, and then by running it on the recent data.
  \item Identify candidate follow-up pull requests for each merged agent PR, and
  measure how often and how soon they occur for each agent (RQ1).
  \item Manually analyse a representative sample of candidate pairs to separate
  corrections from other changes, and classify what was wrong in the agent code
  (RQ2) and how it was corrected (RQ3).
  \item Compare the results across agents and projects and discuss their
  implications for developers, reviewers, and agent builders.
\end{enumerate}

\section{Significance}
\label{sec:significance}

\subsection{Why This Work Matters}
\label{subsec:why-this-work-matters}

AI coding agents are already producing millions of pull requests, and a large
share of them are merged into real projects
\cite{aidev2026hf,watanabe2025agentic}. Every merged agent PR that later needs a
correction creates hidden work for maintainers: someone has to notice the
problem, understand code they did not write, and repair it. If this work is not
measured, the benefits of AI agents may be overestimated.

Knowing \emph{what} goes wrong in merged agent code, and \emph{how} it is
repaired, is useful for several groups. Developers and reviewers can learn which
parts of an agent PR need the most attention before they merge it. Teams can
decide which checks to add to their review process or continuous integration.
Agent builders can use the recurring deficiencies to improve their agents.
Researchers gain an updated, post-merge view of agent code based on recent data,
whereas earlier studies stop in mid-2025.

\subsection{Contributions}
\label{subsec:contributions}

This thesis makes the following contributions:

\begin{enumerate}
  \item \textbf{A review of datasets on AI-generated code changes}, describing
  for each source its size, time period, agents, and how the agent is involved,
  and explaining why AIDev is the main source of agent pull requests.
  \item \textbf{An adaptation of a published follow-up linker to AIDev version
  5}, which no longer contains the task-type and human pull-request tables the
  original linker requires, together with its validation against the original
  implementation and published results.
  \item \textbf{A dataset of candidate follow-up pairs} linking merged agent pull
  requests to later pull requests in Python projects, up to 2026, for further
  research.
  \item \textbf{An empirical analysis of post-merge corrections} of agent code:
  their frequency per agent (RQ1), a classification of the deficiencies they
  address (RQ2), and a description of how they are repaired and by whom (RQ3).
  \item \textbf{A replication package} with the code and data needed to reproduce
  the results.
\end{enumerate}

\section{Definitions}
\label{sec:definitions}

\begin{description}
  \item[AI coding agent] An AI system that receives a development task and
  autonomously edits code, runs commands, and submits the result, typically as a
  pull request (e.g., OpenAI Codex, Devin, Claude Code).
  \item[Pull request (PR)] A request to merge a set of code changes (one or more
  commits) into a repository. It is reviewed and then merged or closed.
  \item[Commit] A single recorded change to the files of a repository, with an
  author and a message. A pull request contains one or more commits.
  \item[Agent pull request] A pull request whose code was generated by an AI
  coding agent, as identified by the AIDev dataset.
  \item[Merged pull request] A pull request whose changes were accepted into the
  main code of the repository.
  \item[Follow-up pull request] A later pull request, merged in the same
  repository within a fixed time window after an agent pull request, that changes
  at least one of the same files.
  \item[Candidate pair] An agent pull request and one of its follow-up pull
  requests. A candidate pair is not yet evidence of a correction.
  \item[Corrective maintenance] A change that removes a fault or deficiency from
  code that is already part of the project \cite{swanson1976dimensions}. In this
  thesis it covers bugs, missing checks, robustness and error-handling problems,
  performance problems, and design or maintainability deficiencies in the agent
  code.
  \item[Linker] The program that builds candidate pairs from repository data
  using the repository, the time window, and the shared files.
  \item[AIDev] A public dataset of pull requests opened by AI coding agents on
  GitHub, with their repositories, commits, reviews, and comments
  \cite{li2026aidev}.
  \item[Observation window] The number of days after the merge of an agent pull
  request in which follow-ups are searched (30 or 60 days in this thesis).
\end{description}
